\begin{theorem}
Let $1<p<\infty$, let $k\ge1$ be an integer and $\alpha\in(0,1]$ with
$k+\alpha$ equal to the right-hand side of (9), as given by
\noderef{KAlphaExp}.  Assume the three quoted bundles (C1), (C2), (C3) of
\noderef{CGradPackage}, \noderef{QCPackage} and \noderef{SmoothOffZeros}.
Let $\Omega\subseteq\mathbb{C}$ be a plane domain, that is open and connected,
and let $u\colon\mathbb{C}\to\mathbb{R}$ be $p$-harmonic on $\Omega$ in the
sense of \noderef{PHarmonicOn}.  Then there is
$\tilde u\colon\mathbb{C}\to\mathbb{R}$ with
\begin{enumerate}
\item[(a)] $\tilde u(z)=u(z)$ for almost every $z$ with $z\in\Omega$;
\item[(b)] $\tilde u\in C^{k,\alpha}_{\mathrm{loc}}(\Omega)$ in the sense of
  \noderef{MemHolderLocOn};
\item[(c)] $\tilde u\in W^{k+2,q}_{\mathrm{loc}}(\Omega)$ in the sense of
  \noderef{MemWLocOn} for every $q$ with $1\le q<2/(2-\alpha)$.
\end{enumerate}

This is the assertion (8) of Theorem 1 of the paper, with the integrability
range (10).  It covers every $1<p<\infty$: Section 6 opens at paper.tex:515
with $p\neq2$, but Theorem 1 claims (8) for all $p$, and only its optimality
clause at paper.tex:95 excludes $p=2$.

Two points about the shape.  First, the conclusion is asserted of a
representative $\tilde u$ rather than of $u$, and it must be: $p$-harmonicity
constrains $u$ only through integrals against test functions and through
membership in $L^{p}_{\mathrm{loc}}$, all unaffected by changing $u$ on a null
set, whereas $C^{k,\alpha}_{\mathrm{loc}}$ requires $C^{k}$-smoothness and so
is destroyed by changing a function at a single point.  Since
$W^{k+2,q}_{\mathrm{loc}}$ is itself invariant under changes on null sets, (c)
holds for $u$ as well as for $\tilde u$.  For the same reason the quoted
bundles (C2) and (C3) are read of a representative of the complex gradient:
quasiregularity and $W^{1,2}_{\mathrm{loc}}$ membership see that gradient only
almost everywhere, so its continuity, the discreteness of its zero set and its
smoothness off that set are statements about its continuous representative, as
the paper's own $f^{-1}(0)=\{z:\nabla u(z)=0\}$ at paper.tex:156 already
indicates.  The proof below introduces that representative once, in Step 2, and
runs on it throughout.  Second, Section 3's normalisation of
the factorization (15) --- that $\chi$ and its inverse may be taken $C^\infty$
off the singular point --- is \emph{not} a hypothesis here: it is derived from
the three bundles by \noderef{NormalizedFactorization}, which are exactly the
ones the note at paper.tex:103 enumerates.  Nothing beyond them is assumed: in
particular the smoothness of the hodograph inverse off the origin, which
Section 3 takes as a standing assumption at paper.tex:183 and which
paper.tex:115 licenses for Theorem 2 only, is not a hypothesis anywhere on this
chain --- it is produced from the expansion (19) by \noderef{Cor1}, as the
paper does at paper.tex:387--388.
\end{theorem}
\begin{proof}
One remark on well-posedness before the argument proper.  The hypothesis above
presents the pair $(k,\alpha)$ as given, whereas Theorem 1 of the paper says at
paper.tex:87--90 that ``the integer $k\ge1$ and the exponent $\alpha\in(0,1]$
are determined by the equation (9)''.  By \noderef{KAlphaWellDefined} there is,
for the given $p$, exactly one pair consisting of an integer $k\ge1$ and a real
number $\alpha$ with $0<\alpha\le1$ such that $k+\alpha$ equals the right-hand
side of (9); so the hypothesis on $(k,\alpha)$ is satisfiable for every
$1<p<\infty$ and determines the pair from $p$, exactly as the paper asserts, and
the statement proved below is therefore not vacuous for any $p$.  This is a
well-posedness remark and not a step in the argument: nothing in the rest of
this proof invokes that lemma, every step below being valid for whichever pair
satisfies the hypothesis.

Write $d=d(p)$ for the exponent of \noderef{DExp} and $d\sigma$ for plane
Lebesgue measure.  By \noderef{KAlphaEqDExpAddOne} and the hypothesis on
$(k,\alpha)$,
\begin{equation}\tag{1}
  d+1=\text{(right-hand side of (9))}=k+\alpha,\qquad\text{so}\qquad
  d=k-1+\alpha .
\end{equation}
In particular $d-m=k+\alpha-(m+1)$ for every $m$, which is the bookkeeping used
below.

\emph{Step 1: the complex gradient.}  By \noderef{PHarmonicOn} there is
$g\colon\mathbb{C}\to\mathbb{C}$ with
$u\in L^{p}_{\mathrm{loc}}(\Omega)$ and $g\in L^{p}_{\mathrm{loc}}(\Omega)$ in
the sense of \noderef{MemLpLocOn}, with
\begin{equation}\tag{2}
  F_0:=\operatorname{Re}g\ \text{the weak }\partial_0\text{-derivative of }u,
  \qquad
  F_1:=\operatorname{Im}g\ \text{the weak }\partial_1\text{-derivative of }u
\end{equation}
on $\Omega$ in the sense of \noderef{IsWeakPartialOn}, and with the weak
$p$-harmonic equation.  Here $F_0$ and $F_1$ are real-valued, and are read as
complex-valued functions wherever the ambient notions require it, as
\noderef{PHarmonicOn} does.  Put
\[
  f=\tfrac12\bar g .
\]
The clause \texttt{complex\_gradient} of \noderef{CGradPackage}, applied with
$p$, $\Omega$, $u$, $g$, the openness of $\Omega$, $1<p$, the
$p$-harmonicity of $u$ and the two clauses (2), gives
\begin{equation}\tag{3}
  f\in W^{1,2}_{\mathrm{loc}}(\Omega),\qquad
  f\ \text{solves (12) on }\Omega,\qquad
  f\ \text{is }K\text{-quasiregular on }\Omega,
\end{equation}
with $K=\max\{p-1,1/(p-1)\}$, in the senses of \noderef{MemWLocOn},
\noderef{SolvesEq12} and \noderef{QuasiregularOn}.  Note $K\ge1$: if $p\ge2$
then $p-1\ge1$, and if $p<2$ then $0<p-1<1$ so $1/(p-1)>1$.

\emph{Step 2: the continuous representative of the gradient.}  First we record
the relation between $g$ and $f$: as $f=\tfrac12\bar g$,
\begin{equation}\tag{4}
  F_0=\operatorname{Re}g=f+\bar f,\qquad
  F_1=\operatorname{Im}g=i\bigl(f-\bar f\bigr),
\end{equation}
the first because $f+\bar f=\tfrac12(\bar g+g)=\operatorname{Re}g$ and the
second because
$i(f-\bar f)=\tfrac{i}{2}(\bar g-g)=\tfrac{i}{2}(-2i\operatorname{Im}g)
=\operatorname{Im}g$.

The clause \texttt{qr\_continuous} of \noderef{QCPackage}, applied with $K$,
$\Omega$, $f$, the openness of $\Omega$ and the quasiregularity in (3),
provides $\hat f\colon\mathbb{C}\to\mathbb{C}$ with
\begin{equation}\tag{4$'$}
  \hat f=f\ \text{almost everywhere on }\Omega,\qquad
  \hat f\ \text{is continuous on }\Omega .
\end{equation}
It does not give continuity of $f$ itself and cannot: by
\noderef{QuasiregularOn} quasiregularity constrains $f$ only through
$W^{1,2}_{\mathrm{loc}}$ membership, which \noderef{MemWLocOn} expresses by
integral identities and indicator $L^2$-conditions, and through an
almost-everywhere inequality between the weak Wirtinger derivatives of $f$,
which by \noderef{IsWeakWirtingerOn} are themselves determined by the class of
$f$ up to null sets; none of this is affected by altering $f$ on a null set.  Put
\[
  \hat F_0=\hat f+\overline{\hat f},\qquad
  \hat F_1=i\bigl(\hat f-\overline{\hat f}\bigr),
\]
which are real-valued and continuous on $\Omega$ by the second clause of
(4$'$), and which by (4) and the first clause of (4$'$) satisfy
$\hat F_j=F_j$ almost everywhere on $\Omega$.  By \noderef{IsWeakPartialOn} the
assertion that a function is the weak $\partial_j$-derivative of $u$ on
$\Omega$ is the identity $\int u\,\partial_j\varphi=-\int F_j\varphi$ over test
functions for $\Omega$, and the right-hand integral is unchanged when $F_j$ is
altered on a null set; so by (2),
\begin{equation}\tag{4$''$}
  \hat F_0\ \text{and}\ \hat F_1\ \text{are the weak }\partial_0\text{- and }
  \partial_1\text{-derivatives of }u\ \text{on }\Omega .
\end{equation}
Finally, the clause \texttt{complex\_gradient\_rep} of
\noderef{CGradPackage}, applied with $p$, $\Omega$, $u$, $g$ and
$\hat f$, the openness of $\Omega$, $1<p$, the $p$-harmonicity of $u$, the two
clauses (2), and the two clauses of (4$'$) --- note that
$\hat f=f=\tfrac12\bar g$ almost everywhere on $\Omega$ is the first of them
--- gives
\begin{equation}\tag{4$'''$}
  \hat f\ \text{is differentiable at almost every point of }\Omega,\qquad
  \hat f\ \text{solves (12) on }\Omega
\end{equation}
in the sense of \noderef{SolvesEq12}.  Of these two clauses it is the first,
the almost-everywhere differentiability of $\hat f$, that is genuinely new:
the second could equally be read off from the second clause of (3), because by
\noderef{SolvesEq12} and \noderef{IsWeakWirtingerOn} solving (12) is a
property of the almost-everywhere class of its argument --- the local
integrability and integration-by-parts clauses see the argument only through
integrals, and its pointwise values enter (12) only inside an ``almost every
$z$'' quantifier --- and $\hat f=f$ almost everywhere on $\Omega$ by (4$'$).
Nothing in (3), by contrast, gives differentiability of any representative at
any point, and that is what the smoothness bundle (C3) consumes.

\emph{Step 3: $u$ is locally integrable.}  Let $\kappa\subseteq\Omega$ be
compact.  By Step 1 the indicator of $\kappa$ times $u$ lies in $L^{p}$, and
$\kappa$ has finite measure, so by H\"older's inequality with exponents $p\ge1$
and $p/(p-1)$ it lies in $L^{1}$.  Hence
$u\in L^{1}_{\mathrm{loc}}(\Omega)$ in the sense of \noderef{MemLpLocOn}.

\emph{Step 4: the representative.}  The hypotheses of
\noderef{WeakGradContinuousRep} hold for $\Omega$, $u$ and
$\hat F_0,\hat F_1$: local integrability of $u$ is Step 3, continuity of
$\hat F_0,\hat F_1$ on $\Omega$ is Step 2, and the weak-derivative clauses are
(4$''$).  It therefore provides
$\tilde u\colon\mathbb{C}\to\mathbb{R}$ with
\begin{equation}\tag{5}
  \tilde u=u\ \text{a.e. on }\Omega,\qquad
  \tilde u\in C^{1}(\Omega),\qquad
  \partial_j\tilde u=\hat F_j\ \text{on }\Omega\ (j=0,1),
\end{equation}
where $\partial_0,\partial_1$ are the partial derivatives of
\noderef{PartialWord} along $[0]$, $[1]$.  The first clause of (5) is
assertion (a).  Combining the third with the definition of $\hat F_j$,
\begin{equation}\tag{6}
  \partial_0\tilde u=\hat f+\overline{\hat f},\qquad
  \partial_1\tilde u=i\bigl(\hat f-\overline{\hat f}\bigr)
  \qquad\text{on }\Omega .
\end{equation}
Throughout the rest of the proof we write $w=\tilde u$, regarded as a
complex-valued function with real values where the ambient notions require it.

\emph{Step 5: the degenerate case $\hat f\equiv0$.}  Suppose $\hat f(z)=0$ for
every $z\in\Omega$.  Then $\hat F_0=\hat F_1=0$ on $\Omega$ by their
definitions, and by (5) both first partial derivatives of $w$ vanish at every
point of $\Omega$.  An
$\mathbb{R}$-linear map $\mathbb{C}\to\mathbb{C}$ is determined by its values
on $1$ and $i$, so $\mathrm{d}w=0$ on $\Omega$; and $w$ is $C^1$ on $\Omega$ by
(5).  A $C^{1}$ function with vanishing derivative on a connected open set is
constant, so $w\equiv c$ on $\Omega$ for some real $c$.

Then (b) and (c) hold.  For (b): $w$ agrees with a constant on the open set
$\Omega$, hence is $C^\infty$ there, so in particular
$\mathrm{ContDiff}^{k}$ on $\Omega$; and the iterated derivative of order
$k\ge1$ of a locally constant function vanishes identically on $\Omega$, so for
any compact $\kappa\subseteq\Omega$ the H\"older condition of
\noderef{MemHolderLocOn} holds with constant $0$.  For (c): let
$1\le q<2/(2-\alpha)$, and take the family $D$ with $D[\,]=w$ and $D(v)=0$ for
every nonempty word $v$.  Each $D(v)$ lies in $L^{q}_{\mathrm{loc}}(\Omega)$,
being either a constant or $0$ and compact sets having finite measure; and for
every word $v$ with $|v|<k+2$ and every $j$, the identity of
\noderef{IsWeakPartialOn} holds: for $v=[\,]$ it reads
$\int w\,\partial_j\varphi=-\int0\cdot\varphi=0$, and
$\int w\,\partial_j\varphi=c\int\partial_j\varphi=0$ because the integral of a
partial derivative of a test function vanishes, as in Step 2 of
\noderef{WeakPartialAcrossPoint}; and for $v$ nonempty both sides are $0$.  So
$w\in W^{k+2,q}_{\mathrm{loc}}(\Omega)$ in the sense of \noderef{MemWLocOn}.
This settles the theorem in this case, and we assume from now on that
$\hat f(z^{*})\neq0$ for some $z^{*}\in\Omega$.

\emph{Step 6: the singular set.}  Put
\[
  Z=\{z\in\Omega:\hat f(z)=0\} .
\]
The clause \texttt{qr\_discrete\_zeros} of \noderef{QCPackage}, applied with
$K$, $\Omega$, the pair $f,\hat f$, the openness and connectedness of
$\Omega$, the quasiregularity in (3), the two clauses of (4$'$) and the point
$z^{*}$, gives for each $z_0\in\Omega$ an $r>0$ with $\hat f(z)\neq0$ for every
$z\in\Omega$ with $z\neq z_0$ and $\mathrm{dist}(z,z_0)<r$.  Read at $z_0\in Z$
this says
\begin{equation}\tag{7}
  \forall z_0\in Z,\ \exists r>0,\ \forall z\in\Omega,\ z\neq z_0,\
  \mathrm{dist}(z,z_0)<r\ \Rightarrow\ z\notin Z .
\end{equation}
Moreover $Z$ meets every compact $\kappa\subseteq\Omega$ in a finite set.
Indeed, $Z$ is closed in $\Omega$, because $\hat f$ is continuous on $\Omega$
by (4$'$); if $Z\cap\kappa$ were infinite it would have an accumulation point
$\zeta\in\kappa\subseteq\Omega$, and $\zeta\in Z$ by closedness, contradicting
(7) at $\zeta$.  So
\begin{equation}\tag{8}
  Z\cap\kappa\ \text{is finite for every compact }\kappa\subseteq\Omega .
\end{equation}

\emph{Step 7: smoothness off the singular set.}  The clause
\texttt{smooth\_off\_zeros} of \noderef{SmoothOffZeros}, applied with $p$,
$\Omega$, the pair $f,\hat f$, the openness of $\Omega$, $1<p$, the first
clause of (3), the two clauses of (4$'$) and the two clauses of (4$'''$), gives
that $\hat f$ is $C^\infty$ on
$\Omega-\{z:\hat f(z)=0\}=\Omega-Z$.  By (6) the two first partial derivatives
of $w$ are the $\mathbb{R}$-linear combinations $\hat f+\overline{\hat f}$ and
$i(\hat f-\overline{\hat f})$ of $\hat f$ and $\overline{\hat f}$, hence are
$C^\infty$ on $\Omega-Z$; and $w$ is $C^{1}$ on $\Omega$ by (5).  A $C^{1}$
function on an open set whose first partial derivatives are $C^\infty$ there is
itself $C^\infty$ there, so
\begin{equation}\tag{9}
  w\ \text{is }C^\infty\ \text{on }\Omega-Z .
\end{equation}

\emph{Step 8: the estimates (51), (52) at each singular point.}  Fix
$\zeta\in Z$, so $\zeta\in\Omega$ and $\hat f(\zeta)=0$.  The hypotheses of
\noderef{NormalizedFactorization} hold with $p$, $K$, the bundles (C2), (C3),
the domain $\Omega$, the pair $f,\hat f$ and the point $\zeta$: the openness
and connectedness of $\Omega$ are hypotheses of the present theorem, the
quasiregularity of $f$ is the third clause of (3), its hypotheses (i) and (ii)
on $\hat f$ are the two clauses of (4$'$) and its (iii) and (iv) are the two
clauses of (4$'''$), the point $z^{*}$ has $\hat f(z^{*})\neq0$ by the case
assumption of Step 5, and $\hat f(\zeta)=0$.  It provides $m\ge1$, $K'\ge1$, an
open $V'$ with $\zeta\in V'\subseteq\Omega$, and $\chi,H$ with: $\chi$
$K'$-quasiconformal on $V'$, $\chi(\zeta)=0$, $\hat f=\chi^{m}$ on $V'$, $\chi$
of class $C^\infty$ on $V'-\{\zeta\}$, $H(\chi(z))=z$ on $V'$, $H$ continuous
on $\chi(V')$, $H(0)=\zeta$, and $H(\xi)\in V'$ with $\chi(H(\xi))=\xi$ for
$\xi\in\chi(V')$.  No smoothness of $H$ is provided and none is needed:
\noderef{Thm1DerivBounds} asks only for its continuity.

The hypotheses of \noderef{Thm1DerivBounds} now hold with $m$, $p$, $K'$, the
bundle (C2), the point $\zeta$, the open set $V'$, the function $w$ in place of
$u$, and $\hat f,\chi,H$ in place of its $f,\chi,H$: (H1) is the restriction to
$V'$ of the second clause of (4$'''$) --- an almost-everywhere statement
restricts to a subset --- together with $\hat f(\zeta)=0$; (H2), (H3), (H4) are
the data just produced; and (H5) is (6) restricted to $V'\subseteq\Omega$.  It therefore provides $\delta_\zeta$ with
$0<\delta_\zeta\le1$ and constants $B^{\zeta}_m$ with
\begin{align}
  \sum_{|v|=m}\bigl|D^{v}\hat f(z)\bigr|
    &\le B^{\zeta}_m|z-\zeta|^{\,d-m},\tag{51$_\zeta$}\\
  \sum_{|v|=m+1}\bigl|D^{v}w(z)\bigr|
    &\le B^{\zeta}_m|z-\zeta|^{\,d-m},\tag{52$_\zeta$}
\end{align}
for $0<|z-\zeta|<\delta_\zeta$ and every $m\ge0$.  Using (1), (52$_\zeta$) reads
in terms of the order $j=m+1$:
\begin{equation}\tag{10}
  \sum_{|v|=j}\bigl|D^{v}w(z)\bigr|\le B^{\zeta}_{j-1}|z-\zeta|^{\,k+\alpha-j}
  \qquad(j\ge1,\ 0<|z-\zeta|<\delta_\zeta).
\end{equation}
Each $B^{\zeta}_{j-1}\ge0$, the left side being a sum of moduli and the power
positive for some admissible $z$.  Shrinking $\delta_\zeta$ we may and do also
assume
\begin{equation}\tag{11}
  \{z:|z-\zeta|<\delta_\zeta\}\subseteq\Omega
  \quad\text{and}\quad
  \{z:|z-\zeta|<\delta_\zeta\}\cap Z=\{\zeta\},
\end{equation}
the first by openness of $\Omega$ and the second by (7).

\emph{Step 9: $w$ is of class $C^{k}$ on $\Omega$.}  The hypotheses of
\noderef{Thm1SmoothAcrossZeros} hold with $\Omega$, $Z$, $k$, $\alpha$ and $w$:
$Z\subseteq\Omega$; the hypothesis that $Z$ meets every compact subset of
$\Omega$ in a finite set is (8); (H1) is the second clause
of (5); (H2) is (9); and (H3) is (10), with $\delta=\delta_\zeta$ and
$B_j=B^{\zeta}_{j-1}$.  Hence
\begin{equation}\tag{12}
  w\ \text{is of class }C^{k}\ \text{on }\Omega .
\end{equation}
This is the first clause of \noderef{MemHolderLocOn}.

\emph{Step 10: from the derivatives of order $k$ to the H\"older condition.}
Write $T(x)=$ the iterated derivative of order $k$ of $w$ within $\Omega$ at
$x$, the object appearing in \noderef{MemHolderLocOn}.  We first record a
comparison with the $D^{v}$.

\emph{(10a) The comparison.}  For $x,y\in\Omega$,
\begin{equation}\tag{13}
  \bigl\|T(x)-T(y)\bigr\|\le\sum_{|v|=k}\bigl|D^{v}w(x)-D^{v}w(y)\bigr| .
\end{equation}
Indeed $T(x)-T(y)$ is a continuous $\mathbb{R}$-multilinear map of $k$
arguments on $\mathbb{C}$, and its norm is the supremum of its values on
$k$-tuples of unit vectors.  Writing a unit vector as $a\cdot1+b\cdot i$ with
$a^{2}+b^{2}=1$, so $|a|,|b|\le1$, and expanding each of the $k$ arguments by
multilinearity, the value on such a tuple is a sum over the $2^{k}$ words $v$
of length $k$ of a coefficient of modulus at most $1$ times the value of
$T(x)-T(y)$ on the corresponding tuple of basis vectors.  That value is the
iterated directional derivative of $w$ along those basis vectors, taken at $x$
and at $y$, which by \noderef{PartialWord} is $D^{v}w(x)-D^{v}w(y)$ for the
word $v$ of the corresponding indices --- the order in which the indices are
read is a convention, and immaterial here because the sum in (13) is over all
words of length $k$.  Taking absolute values and summing gives (13).

\emph{(10b) The three cases.}  Let $\kappa\subseteq\Omega$ be compact; we must
produce $C\ge0$ with $\|T(x)-T(y)\|\le C|x-y|^{\alpha}$ for all
$x,y\in\kappa$.  Since $\Omega$ is open and $\kappa$ compact, fix $r_0>0$ with
$\kappa_{3r_0}:=\{x:\mathrm{dist}(x,\kappa)\le3r_0\}\subseteq\Omega$; this set
is compact.  By (8) the set $Z\cap\kappa_{3r_0}$ is finite.  Choose
$r\in(0,r_0]$ with
\[
  r\le1,\qquad 4r<\delta_\zeta\ \text{ for every }\zeta\in Z\cap\kappa_{3r_0},
  \qquad 4r<\min|\zeta-\zeta'|
\]
over distinct $\zeta,\zeta'\in Z\cap\kappa_{3r_0}$ (the last condition omitted
if there are fewer than two such points).  Let $x,y\in\kappa$.

\emph{Case 1: $|x-y|\ge r$.}  By (12) the map $T$ is continuous on $\Omega$,
hence bounded on the compact $\kappa$, say by $M_T$.  Then
\[
  \|T(x)-T(y)\|\le2M_T=2M_Tr^{-\alpha}r^{\alpha}\le2M_Tr^{-\alpha}|x-y|^{\alpha}.
\]

\emph{Case 2: $|x-y|<r$ and $|x-\zeta|\le2r$ for some $\zeta\in Z$.}  Such a
$\zeta$ lies in $\kappa_{3r_0}$, since $\mathrm{dist}(\zeta,\kappa)\le2r\le3r_0$.
Also $|y-\zeta|\le|y-x|+|x-\zeta|<3r$.  Fix a word $v$ of length $k$ and put
\[
  U=\{\eta:|\eta|<4r\},\qquad F(\eta)=D^{v}w(\eta+\zeta)\quad(\eta\in U).
\]
We verify the hypotheses of \noderef{Lemma2} for $F$ on $U$, in dimension
$n=2$.  The plane $\mathbb{C}$ with its modulus is carried onto
$\mathbb{R}^{2}$ with its Euclidean norm by the $\mathbb{R}$-linear isometry
$\eta\mapsto(\operatorname{Re}\eta,\operatorname{Im}\eta)$, under which
convexity, openness, membership of the origin, continuity, membership of
$C^{1}$, norms of vectors and operator norms of Fréchet derivatives all
correspond; so it is equivalent, and we argue in $\mathbb{C}$.

$U$ is open and convex and contains $0$.  By (11) and $4r<\delta_\zeta$ the
translate $\zeta+U$ is contained in $\Omega$ and meets $Z$ only at $\zeta$.
Hence $F$ is continuous on $U$, by (12); and $F$ is $C^{1}$ on $U-\{0\}$, by
(9).  For hypothesis (i), (10) at $j=k$ gives for $0<|\eta|<4r<\delta_\zeta$
\[
  |F(\eta)|\le\sum_{|v'|=k}\bigl|D^{v'}w(\eta+\zeta)\bigr|
  \le B^{\zeta}_{k-1}|\eta|^{\,k+\alpha-k}=B^{\zeta}_{k-1}|\eta|^{\alpha},
\]
and at $\eta=0$ the same bound holds because the right-hand side tends to $0$
as $\eta\to0$ while $F$ is continuous at $0$, forcing $F(0)=0$.  For
hypothesis (ii), let $0<|\eta|<4r$.  As in Step 9, the operator norm of the
Fréchet derivative of $F$ at $\eta$ is at most the sum of the moduli of its two
first partial derivatives, which are $D^{[0]+\!\!+v}w$ and $D^{[1]+\!\!+v}w$ at
$\eta+\zeta$ --- prefixing a letter applies that partial derivative last, by the
recursion of \noderef{PartialWord} --- and both words have length $k+1$, so by
(10) at $j=k+1$
\[
  \|\mathrm{d}F(\eta)\|\le\sum_{|v'|=k+1}\bigl|D^{v'}w(\eta+\zeta)\bigr|
  \le B^{\zeta}_{k}|\eta|^{\,k+\alpha-(k+1)}=B^{\zeta}_{k}|\eta|^{\alpha-1}.
\]
So \noderef{Lemma2} applies with $M=\max\{B^{\zeta}_{k-1},B^{\zeta}_{k}\}\ge0$
and this $\alpha\in(0,1]$, and yields
\[
  \bigl|D^{v}w(x)-D^{v}w(y)\bigr|=\bigl|F(x-\zeta)-F(y-\zeta)\bigr|
  \le5M|x-y|^{\alpha},
\]
the two arguments lying in $U$ because $|x-\zeta|\le2r<4r$ and
$|y-\zeta|<3r<4r$.  Summing over the $2^{k}$ words $v$ and using (13),
\[
  \|T(x)-T(y)\|\le 2^{k}\cdot5M_{*}|x-y|^{\alpha},\qquad
  M_{*}=\max_{\zeta\in Z\cap\kappa_{3r_0}}
    \max\{B^{\zeta}_{k-1},B^{\zeta}_{k}\},
\]
a maximum over a finite set by (8).

\emph{Case 3: $|x-y|<r$ and $|x-\zeta|>2r$ for every $\zeta\in Z$.}  Then the
ball $\{\xi:|\xi-x|<r\}$ is contained in $\Omega-Z$: it lies in
$\kappa_{3r_0}\subseteq\Omega$, and any $\xi$ in it has
$|\xi-\zeta|\ge|x-\zeta|-|x-\xi|>2r-r=r>0$ for every $\zeta\in Z$.  Put
\[
  S=\bigl\{\xi\in\kappa_{3r_0}:\mathrm{dist}(\xi,Z)\ge r\bigr\},
\]
a compact subset of $\Omega-Z$, and let $L$ be the supremum over $S$ and over
the $2^{k}$ words $v$ of length $k$ of
$\|\mathrm{d}\bigl(D^{v}w\bigr)\|$, finite because each $D^{v}w$ is $C^{1}$ on
the open $\Omega-Z\supseteq S$ by (9), so its derivative is continuous on the
compact $S$.  Every point of the segment from $x$ to $y$ lies in
$\{\xi:|\xi-x|<r\}$, hence in $S$; the segment is convex, so the mean value
inequality gives, for each word $v$ of length $k$,
\[
  \bigl|D^{v}w(x)-D^{v}w(y)\bigr|\le L|x-y|
  =L|x-y|^{\alpha}|x-y|^{1-\alpha}\le L\,r^{1-\alpha}|x-y|^{\alpha},
\]
using $|x-y|<r$ and $1-\alpha\ge0$.  Summing over $v$ and using (13),
$\|T(x)-T(y)\|\le2^{k}Lr^{1-\alpha}|x-y|^{\alpha}$.

Taking for $C$ the maximum of the three constants
$2M_Tr^{-\alpha}$, $2^{k}\cdot5M_{*}$ and $2^{k}Lr^{1-\alpha}$ --- and noting
that the three cases exhaust the possibilities --- gives the H\"older condition
on $\kappa$.  With (12) this proves (b).

\emph{Step 11: the Sobolev clause.}  Let $q$ satisfy $1\le q<2/(2-\alpha)$.
Define a family of functions indexed by words by
\[
  D(v)=D^{\,\mathrm{rev}(v)}w ,
\]
where $\mathrm{rev}(v)$ is $v$ read backwards.  This convention is chosen so
that the coherence required by \noderef{MemWLocOn} is a tautology at the level
of functions: $D([\,])=w$, and for every $v$ and $j$, by the recursion of
\noderef{PartialWord},
\begin{equation}\tag{14}
  D\bigl(v+\!\!+[j]\bigr)=D^{\,j::\mathrm{rev}(v)}w
  =\partial_j\bigl(D^{\,\mathrm{rev}(v)}w\bigr)=\partial_j\bigl(D(v)\bigr),
\end{equation}
so no commutation of mixed partial derivatives is needed anywhere.

\emph{(11a) Local $L^{q}$ bounds.}  Let $v$ be a word with $|v|\le k+2$ and let
$\kappa\subseteq\Omega$ be compact.  If $|v|\le k$ then $D(v)$ is continuous on
$\Omega$ by (12), hence bounded on $\kappa$, hence in $L^{q}$ there, $\kappa$
having finite measure.  Suppose $|v|=j\in\{k+1,k+2\}$.  By (8) the set
$Z\cap\kappa$ is finite; enlarge $\kappa$ to the compact
$\kappa'=\{x:\mathrm{dist}(x,\kappa)\le\varrho\}\subseteq\Omega$ for some
$\varrho>0$, and for each $\zeta\in Z\cap\kappa'$ pick
$\varrho_\zeta\in(0,\min\{\varrho,\delta_\zeta,1\})$.  Off the union of the
balls $\{|z-\zeta|<\varrho_\zeta\}$ the function $D(v)$ is continuous on a
compact subset of $\Omega-Z$ by (9), hence bounded there, so contributes a
finite amount to $\int_{\kappa}\|D(v)\|^{q}$.  On each such ball, (10) gives
$\|D(v)(z)\|\le B^{\zeta}_{j-1}|z-\zeta|^{\,k+\alpha-j}$, and since
$|z-\zeta|<1$ and $k+\alpha-j\ge\alpha-2$ we get
$|z-\zeta|^{\,k+\alpha-j}\le|z-\zeta|^{\alpha-2}$; hence
\[
  \int_{\{|z-\zeta|<\varrho_\zeta\}}\|D(v)\|^{q}\,d\sigma
  \le\bigl(B^{\zeta}_{j-1}\bigr)^{q}
  \int_{\{|z-\zeta|<\varrho_\zeta\}}|z-\zeta|^{(\alpha-2)q}\,d\sigma<\infty
\]
by \noderef{PowerIntegrableNearZero}, applied with this $\alpha$, the radius
$\varrho_\zeta$ and this $q$, after the translation $z\mapsto z-\zeta$, which
preserves Lebesgue measure.  Summing the finitely many contributions,
$D(v)\in L^{q}(\kappa)$.  As $\kappa$ was arbitrary,
$D(v)\in L^{q}_{\mathrm{loc}}(\Omega)$ in the sense of \noderef{MemLpLocOn}.
(The case $q=1$ of this computation will be used in (11b).)

\emph{(11b) The weak-derivative clauses.}  Let $v$ be a word with $|v|<k+2$ and
let $\iota\in\{0,1\}$; we apply \noderef{WeakPartialAcrossPoint} with
$\Omega$, the exceptional set $S=Z$, the index $\iota$, the functions
$G=D(v)$ and $V=D(v+\!\!+[\iota])$, and the exponent $\beta=1-\alpha$, which lies
in $[0,1)$ because $0<\alpha\le1$.  Its hypotheses hold:
$Z\subseteq\Omega$; local finiteness is (8); $G$ is $C^{1}$ on $\Omega-Z$ by
(9); $\partial_\iota G=V$ on $\Omega-Z$ --- indeed everywhere --- by (14); $G$ and
$V$ lie in $L^{1}_{\mathrm{loc}}(\Omega)$ by (11a) with $q=1$, since
$|v|\le k+1$ and $|v+\!\!+[\iota]|\le k+2$; and for the growth bound, let
$\zeta\in Z$ and $0<|z-\zeta|<\min\{\delta_\zeta,1\}$, so that by (10) at the
order $|v|$ --- or trivially if $|v|=0$, $G=w$ being continuous hence bounded
near $\zeta$ ---
\[
  \|G(z)\|\le B^{\zeta}_{|v|-1}\,|z-\zeta|^{\,k+\alpha-|v|}
  \le B^{\zeta}_{|v|-1}\,|z-\zeta|^{\alpha-1}
  =B^{\zeta}_{|v|-1}\,|z-\zeta|^{-\beta},
\]
using $k+\alpha-|v|\ge k+\alpha-(k+1)=\alpha-1$ and $|z-\zeta|<1$.  Hence
\noderef{WeakPartialAcrossPoint} gives that $D(v+\!\!+[\iota])$ is the
distributional $\iota$-th partial derivative of $D(v)$ on $\Omega$, in the sense
of \noderef{IsWeakPartialOn}.

By (11a), (11b) and $D([\,])=w$, the family $D$ witnesses
$w\in W^{k+2,q}_{\mathrm{loc}}(\Omega)$ in the sense of \noderef{MemWLocOn}.
As $q$ was arbitrary in the range $1\le q<2/(2-\alpha)$, this is (c), and with
(a) from Step 4 and (b) from Step 10 the proof is complete.
\end{proof}
